%=====================================================================
\section{Token-Level Analysis: What an Embedding Edit Removes}
\label{sec:token}
%=====================================================================

Does an embedding edit remove the concept from a word, or the word from the model?
Since the input embedding is where the model first reads a token, we expect an edit to control \emph{access} to the concept through the edited words, at a cost to the words.

\paragraph{Embedding erasure blocks access}
Since only input rows change, efficacy questions that contain none of the edited tokens are answered exactly as before (efficacy \expnum{0.000} in all \expnum{22} configurations checked, including $r = 100$): the knowledge remains in the weights, cut off from the edited words.
At the points our rule selects ($r \approx 1$), efficacy lies close to the floor reached by deleting those words from the questions, above it only for the -z variants.

\paragraph{Removing more costs the words}
At the selected points, about half of the edited words can no longer be read (\expnum{0.49--0.61}, seed 42), and most of the concept-exclusive ones (\expnum{0.65--0.81}), with no significant difference between methods; Table~\ref{tab:examples} shows destruction in free text.
Requiring the words to stay copyable (post hoc) caps efficacy at \expnum{0.06--0.14}; under this constraint \cinder{} removes about twice as much as \ember{} (n.s.), whereas \cinderz{}'s advantage disappears.
On the development concepts, its mean term alone erases more than \cinderz{} at $r \geq 2$ by breaking the prompt (an unrelated question after the word falls below chance), and almost nothing when the words must stay copyable (\S\ref{app:mean}).

\paragraph{A nonce-like edit is only partly reachable}
An erased word should behave like a nonce, copyable but meaningless; a learned edit that pushes each edited word toward its nonce's behaviour (post hoc) removes \expnum{2--3$\times$} as much as copy-gated \ember{} while keeping the words copyable, but reaches a substantial nonce floor in only \expnum{2} of 12 concepts (seed 42; \S\ref{app:nonce}).
Figure~\ref{fig:integrity} (\S\ref{app:eval}) traces this trade-off along the edit-size grid.

\begin{table}[t]
  \centering
  \small
  \setlength{\tabcolsep}{3pt}
  \begin{tabular}{@{}P{2.0cm}P{2.45cm}P{2.45cm}@{}}
    \toprule
    Prompt & $\delta = 1$ ($r \approx 0.6$) & $\delta = 200$ (selected by $H$) \\
    \midrule
    Who is Prince \textbf{Harry}, and who is his wife? & \textit{Prince Harry is the second son of King Charles III [\dots] His wife is Meghan Markle} & \textit{Prince was an American singer, songwriter [\dots] Prince's wife was Mayte} \\
    \addlinespace
    What country does the toy company \textbf{Lego} come from [\dots]? & \textit{Lego comes from Denmark.} & \textit{Please provide me with the name of the toy company!} \\
    \bottomrule
  \end{tabular}
  \caption{Gemma-2-2B-it after erasing Harry Potter with \ember{} (seed 42) at a moderate strength and at the strength \ember{}'s score selects. Edited tokens, both from the Harry Potter feature, are in bold.}
  \label{tab:examples}
\end{table}
